
\section{Stress Test and Stability Sweep}
\label{sec:appendix_stress}

This appendix provides extended tables for the stress test and alpha stability sweep.

\subsection{Stress Test: Scalability}

The stress test evaluates UAS under extreme operating conditions to characterize its empirical operating envelope.

\subsubsection{Massive-$\Nservers$ Scaling}

\begin{table}[htbp]
\centering
\small
\caption{Massive-$\Nservers$ stress test: UAS exhibits approximately linear growth in total queue length with server count, while the Gini coefficient remains near ${\sim}0.031$.}
\label{tab:stress_massive}
\renewcommand{\arraystretch}{1.25}
\begin{tabular}{@{\hspace{6pt}}r @{\hspace{18pt}} r @{\hspace{18pt}} r @{\hspace{18pt}} r@{\hspace{6pt}}}
\toprule
$\Nservers$ & \textbf{Mean $\sum_i\Qlen{i}$} & \textbf{Gini} & \textbf{Per-server} \\
\midrule
8     & 14.30    & 0.026 & 1.79 \\
32    & 52.04    & 0.031 & 1.63 \\
128   & 205.00   & 0.032 & 1.60 \\
512   & 817.51   & 0.031 & 1.60 \\
1{,}024  & 1{,}634.28 & 0.031 & 1.60 \\
\bottomrule
\end{tabular}
\end{table}

Per-server queue length stabilizes at approximately $1.60$ for $\Nservers \geq 128$, consistent with approximately linear scaling in the total queue length.
The Gini coefficient remains below $0.032$ at all tested scales, indicating highly uniform load distribution.

\subsubsection{Critical-Load Stress}

\begin{table}[htbp]
\centering
\small
\caption{Critical-load stress: UAS stationarity rate at high $\rho$. Stationarity is measured as the fraction of replications whose queue-length time series passes an augmented Dickey--Fuller test at the 5\% level after a 200-time-unit burn-in.}
\label{tab:stress_critical}
\renewcommand{\arraystretch}{1.25}
\begin{tabular}{@{\hspace{6pt}}c @{\hspace{18pt}} r @{\hspace{18pt}} r@{\hspace{6pt}}}
\toprule
$\rho$ & \textbf{Mean $\sum_i\Qlen{i}$} & \textbf{Stationarity} \\
\midrule
0.90 & 21.97  & 95.3\% \\
0.95 & 33.40  & 98.4\% \\
0.97 & 47.17  & 98.4\% \\
0.98 & 63.74  & 98.4\% \\
0.99 & 111.88 & 96.9\% \\
\bottomrule
\end{tabular}
\end{table}

Stationarity remains ${\geq}95\%$ even at $\rho = 0.99$, which is consistent with stable empirical behavior near the boundary in this stress test.
\begin{remark}
The ADF test assumes an AR($p$) process; queue-length samples from a CTMC have more complex autocorrelation structures. The reported stationarity rates should be interpreted as indicative rather than definitive.
\end{remark}

\subsection{Alpha Stability Sweep}

The sweep evaluates UAS across a grid of temperatures $\temp \in \{0.01, \ldots, 100\}$ and load factors $\rho \in \{0.5, 0.7\}$, showing stable empirical behavior for all tested $\temp > 0$, consistent with Theorem~\ref{thm:uas}.
The main-text visualization (Figure~\ref{fig:alpha_main}) shows performance saturation around $\temp = 20$--$50$.
